\documentclass[12pt]{article}

\usepackage[utf8]{inputenc}
\usepackage{amsmath, amssymb, amsthm}
\usepackage{geometry}
\usepackage{enumitem}
\usepackage{hyperref}

\geometry{margin=1in}

\title{Linear Algebra -- Study Notes}
\author{}
\date{}

\begin{document}
\maketitle

\section{Inner Product Properties (Derived from Axioms)}

The inner product axioms are:
\begin{enumerate}
    \item \textbf{Positive-definiteness:} $\langle v, v \rangle \ge 0$, with $\langle v, v \rangle = 0 \iff v = 0$.
    \item \textbf{Conjugate symmetry:} $\langle u, v \rangle = \overline{\langle v, u \rangle}$.
    \item \textbf{Linearity in the first variable:} $\langle ru + v, w \rangle = r\langle u, w \rangle + \langle v, w \rangle$.
\end{enumerate}

\subsection{Property 1: Additivity in the second variable}

\textbf{Claim:} $\langle u, v+w \rangle = \langle u, v \rangle + \langle u, w \rangle$.

\begin{proof}
By conjugate symmetry:
$\langle u, v+w \rangle = \overline{\langle v+w, u \rangle}$.

By linearity in the first variable:
$= \overline{\langle v, u \rangle + \langle w, u \rangle}$.

Conjugate distributes over addition:
$= \overline{\langle v, u \rangle} + \overline{\langle w, u \rangle}$.

By conjugate symmetry on each term:
$= \langle u, v \rangle + \langle u, w \rangle$.
\end{proof}

\subsection{Property 2: Conjugate-linearity in the second variable}

\textbf{Claim:} $\langle u, rv \rangle = \bar{r}\,\langle u, v \rangle$.

\begin{proof}
By conjugate symmetry:
$\langle u, rv \rangle = \overline{\langle rv, u \rangle}$.

By linearity in the first variable:
$= \overline{r\,\langle v, u \rangle}$.

Conjugate of a product ($\overline{r \cdot z} = \bar{r} \cdot \bar{z}$):
$= \bar{r} \cdot \overline{\langle v, u \rangle}$.

By conjugate symmetry:
$= \bar{r}\,\langle u, v \rangle$.
\end{proof}

\subsection{Property 3: Inner product with zero}

\textbf{Claim:} $\langle u, 0 \rangle = 0$.

\begin{proof}
Write $0 = 0 \cdot v$ for any $v$. By Property~2:
$\langle u, 0 \cdot v \rangle = \bar{0} \cdot \langle u, v \rangle = 0$.

Similarly, $\langle 0, u \rangle = 0$ by conjugate symmetry:
$\langle 0, u \rangle = \overline{\langle u, 0 \rangle} = \overline{0} = 0$.
\end{proof}

\subsection{Property 4: Cancellation}

\textbf{Claim:} If $\langle u, v \rangle = \langle u, w \rangle$ for all $u$, then $v = w$.

\begin{proof}
Subtracting: $\langle u, v \rangle - \langle u, w \rangle = 0$ for all $u$.

We rewrite the left side as $\langle u, v - w \rangle$ by expanding:
\[
\langle u, v + (-1)w \rangle
= \langle u, v \rangle + \langle u, (-1)w \rangle
\quad \text{(Property~1: additivity in second slot)}
\]
\[
= \langle u, v \rangle + \overline{(-1)}\,\langle u, w \rangle
= \langle u, v \rangle - \langle u, w \rangle
\quad \text{(Property~2 with $r = -1$: $\bar{r} = -1$)}
\]
So $\langle u, v - w \rangle = 0$ for all $u$.

Choose $u = v - w$:
$\langle v-w, v-w \rangle = 0$.

By positive-definiteness: $v - w = 0$, so $v = w$.
\end{proof}

\subsection{Property 5: $\langle T(v), w \rangle = 0$ for all $v, w$ implies $T = 0$ (over $\mathbb{C}$)}

\textbf{Claim:} If $T: V \to V$ is linear and $\langle T(v), w \rangle = 0$ for all $v, w \in V$, then $T = 0$.

\begin{proof}
Set $v = ru + w$ for arbitrary $r \in \mathbb{C}$, $u, w \in V$. By hypothesis:
\[
0 = \langle T(ru + w), w \rangle.
\]

Linearity of $T$ and linearity in the first slot:
\[
0 = r\,\langle T(u), w \rangle + \langle T(w), w \rangle.
\]

Fix $u$ and $w$. The term $\langle T(w), w \rangle$ is a fixed complex number, but $r$ is free to be any complex scalar. The only way $r \cdot c = \text{constant}$ holds for every $r \in \mathbb{C}$ is if $c = 0$. Therefore:
\[
\langle T(u), w \rangle = 0 \quad \text{for all } u, w.
\]

Pick $w = T(u)$:
$\langle T(u), T(u) \rangle = 0$, so $T(u) = 0$ by positive-definiteness.

Since $u$ was arbitrary, $T = 0$.
\end{proof}

\section{Cauchy--Schwarz Inequality}

\textbf{Theorem.} For all $u, v$ in an inner product space:
\[
|\langle u, v \rangle| \le \|u\|\|v\|.
\]

\begin{proof}
If $v = 0$, both sides are $0$. Assume $v \neq 0$.

For any $\lambda \in \mathbb{C}$, positive-definiteness gives:
\[
0 \le \|u - \lambda v\|^2 = \langle u - \lambda v,\; u - \lambda v \rangle.
\]
Expanding by linearity in the first slot and conjugate-linearity in the second:
\[
= \|u\|^2 - \lambda\overline{\langle u,v\rangle} - \bar{\lambda}\langle u,v\rangle + |\lambda|^2\|v\|^2.
\]
(Using: $\langle \lambda v, u \rangle = \lambda\langle v, u\rangle = \lambda\overline{\langle u,v\rangle}$ and $\langle u, \lambda v\rangle = \bar{\lambda}\langle u,v\rangle$.)

Choose $\lambda = \dfrac{\langle u,v\rangle}{\|v\|^2}$. Then:
\begin{itemize}
    \item $\bar{\lambda} = \dfrac{\overline{\langle u,v\rangle}}{\|v\|^2}$
    \item $|\lambda|^2 = \dfrac{|\langle u,v\rangle|^2}{\|v\|^4}$
    \item $\lambda\overline{\langle u,v\rangle} = \dfrac{\langle u,v\rangle \cdot \overline{\langle u,v\rangle}}{\|v\|^2} = \dfrac{|\langle u,v\rangle|^2}{\|v\|^2}$
    \item $\bar{\lambda}\langle u,v\rangle = \dfrac{\overline{\langle u,v\rangle} \cdot \langle u,v\rangle}{\|v\|^2} = \dfrac{|\langle u,v\rangle|^2}{\|v\|^2}$
    \item $|\lambda|^2\|v\|^2 = \dfrac{|\langle u,v\rangle|^2}{\|v\|^2}$
\end{itemize}
So all three trailing terms equal $\dfrac{|\langle u,v\rangle|^2}{\|v\|^2}$, giving:
\[
0 \le \|u\|^2 - \frac{|\langle u,v\rangle|^2}{\|v\|^2}.
\]
Rearranging: $|\langle u,v\rangle|^2 \le \|u\|^2\|v\|^2$. Taking square roots:
\[
|\langle u,v\rangle| \le \|u\|\|v\|. \qedhere
\]
\end{proof}

\textbf{Equality} holds iff $u = \lambda v$ for some scalar $\lambda$ (i.e., $u$ and $v$ are linearly dependent).

\subsection{Why choose $\lambda = \langle u,v\rangle / \|v\|^2$?}

This is the value that minimizes $\|u - \lambda v\|^2$ over $\lambda$. Geometrically, $\lambda v$ is the orthogonal projection of $u$ onto $v$, and the inequality says the residual $u - \lambda v$ has non-negative squared norm.

\section{Triangle Inequality (for the induced norm)}

\textbf{Theorem.} $\|u + v\| \le \|u\| + \|v\|$.

\begin{proof}
Expand:
\begin{align*}
\|u+v\|^2 &= \langle u+v,\; u+v\rangle \\
&= \|u\|^2 + \langle u,v\rangle + \langle v,u\rangle + \|v\|^2 \\
&= \|u\|^2 + \|v\|^2 + 2\operatorname{Re}\langle u,v\rangle.
\end{align*}
The last step uses conjugate symmetry: $\langle u,v\rangle + \langle v,u\rangle = \langle u,v\rangle + \overline{\langle u,v\rangle} = 2\operatorname{Re}\langle u,v\rangle$.

Now bound using Cauchy--Schwarz:
\[
\operatorname{Re}\langle u,v\rangle \le |\langle u,v\rangle| \le \|u\|\|v\|.
\]
Therefore:
\[
\|u+v\|^2 \le \|u\|^2 + 2\|u\|\|v\| + \|v\|^2 = \bigl(\|u\| + \|v\|\bigr)^2.
\]
Taking square roots: $\|u+v\| \le \|u\| + \|v\|$.
\end{proof}

\section{Polarization Identities}

The polarization identities recover the inner product from the norm alone.

\subsection{Real case}

\textbf{Theorem.} Over $\mathbb{R}$: $\displaystyle\langle x,y\rangle = \frac{1}{4}\bigl(\|x+y\|^2 - \|x-y\|^2\bigr)$.

\begin{proof}
From expanding via the inner product:
\begin{align*}
\|x+y\|^2 &= \|x\|^2 + \|y\|^2 + \langle x,y\rangle + \langle y,x\rangle, \\
\|x-y\|^2 &= \|x\|^2 + \|y\|^2 - \langle x,y\rangle - \langle y,x\rangle.
\end{align*}
Subtracting:
\[
\|x+y\|^2 - \|x-y\|^2 = 2\langle x,y\rangle + 2\langle y,x\rangle.
\]
Over $\mathbb{R}$, $\langle y,x\rangle = \overline{\langle x,y\rangle} = \langle x,y\rangle$ (since the inner product is real-valued), so the right side is $4\langle x,y\rangle$.
\end{proof}

\subsection{Complex case}

\textbf{Theorem.} Over $\mathbb{C}$:
\[
\langle x,y\rangle = \frac{1}{4}\bigl(\|x+y\|^2 - \|x-y\|^2 + i\|x+iy\|^2 - i\|x-iy\|^2\bigr).
\]

\begin{proof}
Over $\mathbb{C}$, $\langle x,y\rangle + \overline{\langle x,y\rangle} = 2\operatorname{Re}\langle x,y\rangle$, so the subtraction from the real case only gives:
\[
\|x+y\|^2 - \|x-y\|^2 = 4\operatorname{Re}\langle x,y\rangle.
\]
To recover $\operatorname{Im}\langle x,y\rangle$, replace $y$ with $iy$. Using $\langle x, iy\rangle = \bar{i}\langle x,y\rangle = -i\langle x,y\rangle$ and $\langle iy, x\rangle = i\langle y,x\rangle = i\overline{\langle x,y\rangle}$:
\begin{align*}
\|x+iy\|^2 &= \|x\|^2 + \|y\|^2 - i\langle x,y\rangle + i\overline{\langle x,y\rangle}, \\
\|x-iy\|^2 &= \|x\|^2 + \|y\|^2 + i\langle x,y\rangle - i\overline{\langle x,y\rangle}.
\end{align*}
Subtracting:
\[
\|x+iy\|^2 - \|x-iy\|^2 = -2i\langle x,y\rangle + 2i\overline{\langle x,y\rangle}.
\]
Writing $\langle x,y\rangle = a+bi$:
\[
-2i(a+bi) + 2i(a-bi) = -2ia + 2b + 2ia + 2b = 4b = 4\operatorname{Im}\langle x,y\rangle.
\]
Combining:
\begin{align*}
\tfrac{1}{4}\bigl(\|x+y\|^2 - \|x-y\|^2\bigr) &= \operatorname{Re}\langle x,y\rangle, \\
\tfrac{1}{4}\bigl(\|x+iy\|^2 - \|x-iy\|^2\bigr) &= \operatorname{Im}\langle x,y\rangle.
\end{align*}
Therefore:
\[
\frac{1}{4}\bigl(\|x+y\|^2 - \|x-y\|^2 + i\|x+iy\|^2 - i\|x-iy\|^2\bigr) = \operatorname{Re}\langle x,y\rangle + i\operatorname{Im}\langle x,y\rangle = \langle x,y\rangle. \qedhere
\]
\end{proof}

\end{document}
